% Literature review, aligned with the v3 individual-large-trade research design.
% Verified references use natbib; remaining legacy citations await metadata verification.

\section{Literature Review}
\label{sec:literature-review}

\subsection{Introduction}
\label{lr:introduction}

Prediction markets are designed to aggregate dispersed information into continuously updated contract prices, although reading those prices as probabilities requires assumptions about preferences, wealth and market structure \citep{hayek1945knowledge,wolferszitzewitz2004prediction,manski2006interpreting}. Experimental and field evidence from political markets and the Iowa Electronic Markets generally supports their forecasting value, while also documenting setting-specific bias and miscalibration \citep{forsythe1992anatomy,berg2008dozen,arrow2008promise}. Polymarket combines an off-chain order-matching service with blockchain settlement and publicly observable transaction records. This produces unusually rich trade data but does not reveal beneficial ownership, private information, order cancellations throughout the historical sample, or trading intent. The defensible question is therefore whether large individual trades are conditionally associated with short-horizon price adjustment, persistence and reversal, rather than whether they cause price changes or constitute misconduct. This review is organised around the primary hypotheses tested in Section~\ref{sec:results}: signed flow and directional price change as a baseline check (H1); large- versus ordinary-trade flow and its extensive-/intensive-margin decomposition, the dissertation's primary contribution (H2); persistence (H3a) and forecast improvement (H3b); and timing and match-phase heterogeneity, posed as a research question rather than a directional hypothesis (RQ4). It also provides the theoretical and empirical grounding for three supplementary market-integrity analyses reported alongside them, concentration and update incidence (MI1), concentrated pressure and reversal (MI2), and out-of-sample reversal prediction (MI3), and for the dissertation's separate past-only high-activity-wallet and whale-wallet constructs.

\subsection{Market microstructure and the price impact of large trades}
\label{lr:microstructure}

The theoretical starting point for this dissertation is the market microstructure literature on how trade size relates to information content and price impact. Informed traders generate systematic, persistent price impact, and market makers widen spreads to protect themselves against adverse selection from better-informed counterparties \citep{kyle1985continuous,glostenmilgrom1985bidask}. This logic extends directly to trade size: larger trades are more likely to originate from informed participants and should therefore carry more information \citep{easleyohara1987tradesize}. A trade's price impact can be decomposed into a permanent, information-driven component and a temporary, liquidity-driven component \citep{hasbrouck1991information,madhavan1997prices}, a decomposition this dissertation adapts to compare the short-horizon price association of large versus ordinary trades (H2) and to distinguish persistence from reversal (H3a). Empirical institutional-order evidence further shows that market impact can be estimated with nonlinear power laws; \citet{almgren2005direct} reject a square-root temporary-impact specification in favour of an exponent near three-fifths in their sample. This is a possible functional-form robustness benchmark, not a justification for imposing that elasticity on Polymarket.

A second strand of this literature concerns how automated and high-frequency activity changes the interpretation of large or clustered order flow. Large institutional trades can generate heavy-tailed price movements \citep{gabaix2006institutional}, a theoretical analogue for the heavy-tailed distribution of individual trade sizes that motivates this dissertation's within-market P99 large-trade threshold. Algorithmic and high-frequency trading can improve liquidity and price discovery \citep{hendershott2011algorithmic,brogaard2014hft}, although high-frequency market makers nonetheless concentrate activity in ways that can blur the line between liquidity provision and price-moving behaviour \citep{menkveld2013marketmakers}. This distinction motivates classifying large trades by their market-specific relative size rather than by inferred trader identity or role: the P99 threshold used in Section~\ref{sec:research-design} is a relative-size rule, not a wallet, ownership or strategy classification. Flow-toxicity and order-imbalance measures, originally developed for equities, supply a template for the signed order-flow measures this dissertation constructs from Polymarket's own trade data \citep{easley2012flowtoxicity,easley2016discerning}, without treating flow toxicity itself as direct evidence of informed trading.

Persistent order flow need not mean that each trade is an independent information event: it may arise because a larger intended position is split into a sequence of smaller executions. Empirical evidence documents this order-splitting mechanism in both aggregate equity flow and individual trading accounts \citep{toth2015orderflow,sato2023ordersplitting}. This is one reason the dissertation treats the individual P99 trade as an observable execution-size construct rather than equating it with a complete latent order or with the sophistication of its wallet, and it motivates testing same-direction continuation among high-activity wallets as a separate, secondary construct rather than as evidence bearing on the large-trade estimand itself.

\subsection{Informed trading, efficiency and trader skill}
\label{lr:informed-trading}

A persistent puzzle in this literature is why informed trading and concentrated profits can coexist with broadly efficient prices. If prices were perfectly informative, no one would have an incentive to gather information, so some profitable informed trading must survive even in efficient markets \citep{grossmanstiglitz1980impossibility}. The classic efficiency framework provides the benchmark against which Polymarket's resolution accuracy can be tested \citep{fama1970efficient}. Evidence that frequent trading and behavioural biases reduce ordinary investors' performance provides a counterpoint to concentrated Polymarket wallet profits \citep{barberodean2000hazardous,odean1998disposition,odean1999toomuch}.

To separate skill from luck at the wallet level, this dissertation also draws on the broader trader-performance literature, although wallet-level profit and loss is not itself estimated here. Persistent skill exists only among a small subset of day traders \citep{barber2014speculatorskill}, trader characteristics predict commodity-futures profits \citep{dewally2013traderprofits}, and informed-insider profits have been quantified in dollar terms \citep{czirakigider2021insiderprofits}. Trading in lottery-like or long-shot positions can also reflect gambling preferences rather than information \citep{kumar2009gambles,dorn2015gambling}. These studies motivate treating some large trades on low-probability Polymarket contracts as potentially belief-driven, non-informational demand; the trade-level data used here cannot directly adjudicate between gambling preference and superior information.

\subsection{Prediction markets and information aggregation}
\label{lr:prediction-markets}

Beyond general market efficiency, a dedicated literature evaluates prediction markets specifically as forecasting instruments. The theoretical and policy case builds on the idea that prices coordinate dispersed knowledge \citep{hayek1945knowledge,wolferszitzewitz2004prediction,arrow2008promise}. Evidence from the Iowa Electronic Markets documents election-forecasting performance \citep{forsythe1992anatomy,berg2008dozen}, while historical presidential betting markets were broadly informative even when manipulation was attempted \citep{rhodestrumpf2004historical}.

Sports markets provide the closest substantive comparison for the present sample. World Cup market evidence tests efficiency in the same tournament setting \citep{gil2007worldcup}, while football studies show that unanticipated match events and evolving match dynamics can generate overreaction, underreaction and rapid repricing \citep{choihui2014surprise,angelini2022inplay,michels2023matchdynamics}. This evidence supports separating pre-kickoff from post-kickoff observations: after kickoff, public match events change the information-generating process rather than merely shortening the remaining horizon. Statistical work on detecting match fixing demonstrates the value of screening market patterns while distinguishing a statistical alert from proof of misconduct \citep{forrestmchale2019matchfixing}; this alert-versus-proof distinction is the direct precedent for the dissertation's unsupervised screening and contextual case audit, which ranks unusual episodes for review but never labels a flagged or reviewed episode as manipulation.

This optimism is tempered by evidence of systematic mispricing. Prediction market prices need not equal traders' mean beliefs, so equating price with probability requires care \citep{manski2006interpreting}. Betting markets exhibit a persistent favourite-long-shot bias, in which low-probability outcomes are systematically overpriced \citep{snowbergwolfers2010favourite,thalerziemba1988betting}, a pattern this dissertation examines through the timing and phase heterogeneity posed as RQ4 rather than through wallet-level activity. Prediction market prices can further be used as an event-study tool to isolate the market-implied effect of political events \citep{snowberg2007partisan}, an approach this dissertation adapts to measure the short-horizon price association of large-trade-driven order flow (H1 and H2).

\subsection{Market manipulation and market integrity}
\label{lr:manipulation}

Prices can be manipulated through ordinary trading even without the dissemination of false information, simply because other traders cannot perfectly distinguish informed from uninformed large orders \citep{allengale1992manipulation}; this is the theoretical foundation for trade-based manipulation. Manipulated prices can further distort real economic decisions that rely on those prices \citep{goldsteinguembel2008manipulation}, which provides the broader motivation for the supplementary market-integrity analysis, beyond any individual trader's own trading profit. Cash-settled contracts are particularly vulnerable to manipulation incentives concentrated around settlement \citep{kumarseppi1992cashsettlement}, a finding directly relevant to Polymarket's binary, cash-settled markets near resolution. Participants who can influence the underlying outcome face heightened manipulation incentives \citep{ottavianisorensen2007outcome}, although attempted manipulation can also attract informed traders who profit from correcting the resulting mispricing \citep{hanson2006aggregation}, so manipulation and informed correction are not always easy to distinguish empirically, a tension this dissertation's H3a persistence/reversal test addresses descriptively rather than conclusively.

Empirically, manipulation effects on racetrack betting markets tend to be limited and often arbitraged away \citep{camerer1998racetrack}, while one of the earliest large-scale empirical signatures of manipulation in equity markets, namely abnormal price run-ups, volume spikes and subsequent reversals around manipulated stocks, has been documented directly \citep{aggarwalwu2006manipulation}. This price-volume-reversal signature motivates both the dissertation's H3a test of whether large-trade-associated price movements persist from minute five to minute fifteen, which is more consistent with information, or reverse, which is more consistent with temporary price pressure, and the supplementary MI2 test of whether concentrated directional pressure specifically predicts subsequent reversal; neither pattern is treated as direct proof of manipulation or informed trading.

\subsection{Crypto markets and Polymarket-specific evidence}
\label{lr:crypto}

Because Polymarket settles on-chain, this dissertation also situates the research within the cryptocurrency market-quality literature. Cryptocurrency markets exhibit persistent frictions and a risk-return profile distinct from conventional equities \citep{makarovschoar2020crypto,liutsyvinski2021crypto}, which justifies treating Polymarket as economically separate from traditional financial markets rather than assuming equity-market intuitions transfer directly. Evidence consistent with coordinated price support has been documented in Bitcoin markets \citep{griffinshams2020untethered}, while cryptocurrency pump-and-dump schemes display abnormal volume, price jumps and reversals \citep{li2025pumpdump}. These are comparators for the dissertation's H3a/H3b design; the present study does not screen for coordinated multi-wallet trading.

Most directly relevant is a fast-growing body of 2025--2026 Polymarket working papers. Their recent and evolving status requires version-specific interpretation rather than treating them as settled peer-reviewed evidence. \citet{akey2026winslosses} document highly concentrated wallet-level gains and a positive association between maker-volume share and profitability, but find only modest month-to-month persistence and caution that continuation in trading may select the observed sample. \citet{dellavedova2026execution} separate directional accuracy from execution quality using terminal payoffs and alternative price benchmarks. In a large resolved-trade sample, the paper reports that the two dimensions are nearly independent and that execution differences account for the principal profitability contrast across its behavioural wallet types. This finding cautions against inferring forecasting skill from either profit or a large short-horizon price association. \citet{gomezcram2026minority} address skill more directly: they compare each account's realised PnL with an event-level sign-randomisation distribution, define groups on randomly selected training events, and evaluate price discovery on held-out events. Their design demonstrates that skill requires a dedicated outcome-aware test; high activity or trade size alone is not a skill measure.

\citet{ng2026pricediscovery} provide the closest audited design comparison. They construct prices and returns on a fixed 30-minute grid and classify a transaction as large when its volume exceeds the 99th percentile of the relevant platform-day distribution. They then aggregate signed buy and sell volume into a normalised order-imbalance measure, reported as net order imbalance (NOI) or order imbalance (OIB), and find that lagged large-trade NOI/OIB predicts next-interval returns in their two-contract, cross-platform election sample. Their unit labelled ``whale'' is therefore an individual large transaction, not a persistent wallet. This supports a percentile-based size classification and a directional-flow mechanism, but NOI/OIB is not interchangeable with this dissertation's main estimand, because the two ask different questions: their model tests whether interval-level directional imbalance predicts later returns, whereas H2 tests whether large- and ordinary-trade flow have different conditional price associations across market-minutes. NOI/OIB is accordingly treated as a candidate mechanism and robustness comparator rather than a replacement for the large-versus-ordinary contrast that defines H2.

\citet{cheongtamayo2026concentrated} provide the closest wallet-by-trade-size comparison. In earnings prediction markets, they rank wallets within each market by the absolute value of their final volume-weighted net position, label the top one, three or five wallets as whales, and separately classify trades above the within-market 75th percentile of USDC value as large. Their price-response regression interacts these two indicators and uses market fixed effects with wallet- and market-clustered errors. The reported interaction is positive, especially in lower-volume markets. This factorial distinction supports analysing wallet scale and transaction size separately, but the ex-post final-position whale rank, P75 cutoff and consecutive-trade absolute log-odds response differ from this dissertation's past-only wallet measures, P99 threshold and fixed-horizon directional outcome.

\citet{tsangyang2026anatomy} supply a complementary Kyle-style analysis of the 2024 presidential market. They infer direction using the tick rule, aggregate trades hourly, transform prices to log odds and estimate a no-intercept rolling 30-day order-flow coefficient for Trump YES. The reported decline in lambda as activity broadens is consistent with a maturing, deeper market, but overlapping windows, inferred signs and limited rolling-series inference make this a descriptive liquidity result rather than a causal manipulation test. Because a Kyle-style coefficient summarises average price sensitivity to signed flow over an aggregation window, it is relevant to this dissertation only as an auxiliary liquidity and price-sensitivity specification reported alongside the main estimates; it is not an automatic replacement for the large-versus-ordinary flow regression that tests H2. \citet{reichenbach2026decentralized} provide a broader descriptive benchmark based on more than 478 million reconstructed trades and a manually matched tennis sample in which Polymarket probabilities are compared with bookmaker odds. Their sports evidence motivates external-validity comparison, but they do not estimate large-trade price response. An unpublished order-book manuscript further reports that feed-inferred and on-chain trade directions agree only imperfectly and that direction-dependent Kyle estimates can change sign \citep{dubach2026microstructure}. This reinforces the need to document direction provenance. Other studies on aggressor-direction reconstruction, copy trading or isolated oracle disputes remain complementary data and integrity references. This hierarchy is important because the dissertation's canonical sample comes from the Polymarket Data API and does not rely on a third-party database to add or remove observations.

\subsection{Synthesis}
\label{lr:synthesis}

Taken together, this literature establishes three things. First, microstructure theory predicts that large trades on a CLOB-based venue should generate a measurable short-horizon price association, while persistence and reversal can help distinguish lasting adjustment from temporary pressure \citep{kyle1985continuous,hasbrouck1991information,madhavan1997prices}; this motivates H1, H2 and H3a and, through wallet-level concentration, the supplementary MI1 and MI2 tests, without identifying causality. Second, prediction-market and manipulation research shows why forecast accuracy, price response and misconduct are non-equivalent outcomes \citep{manski2006interpreting,hanson2006aggregation,aggarwalwu2006manipulation}; this distinction is carried through the dissertation's separate treatment of association (H1, H2, H3a, H3b and RQ4, MI1--MI2), out-of-sample prediction (MI3), and unresolved anomaly review. Third, the emerging Polymarket literature shows that apparently similar ``large participant'' results depend strongly on the construct being measured. Final-position rank, behavioural wallet type, randomisation-based skill, individual trade size and interval order imbalance are empirically distinct \citep{dellavedova2026execution,gomezcram2026minority,cheongtamayo2026concentrated,ng2026pricediscovery}. This distinction is the basis for separating the dissertation's primary large-trade estimand from its secondary past-only wallet analyses.

\subsection{Research gap}
\label{lr:research-gap}

The literature provides direct precedent for percentile-based large-trade classification and for testing whether aggregated large-trade imbalance predicts later returns, but evidence remains concentrated in political or financial-event contracts and short single-event samples. It does not establish whether the large-versus-ordinary flow association generalises across a broad set of related sports contracts with verified kickoff times and distinct pre- and post-kickoff information environments. This is the setting-specific and design-specific gap addressed here. Using 6,725,732 canonical trades across 360 World Cup markets, the dissertation applies a market-specific P99 trade-value threshold, directly tests the large-minus-ordinary coefficient contrast, separates update incidence from conditional magnitude, and evaluates persistence, Brier improvement and sports-phase heterogeneity, alongside a supplementary market-integrity analysis of wallet concentration, reversal and out-of-sample predictability. The unit of primary classification is the individual trade; wallet activity is analysed separately, and all conclusions, including those of the supplementary analysis, remain associative rather than evidence of manipulation.
